\section{Introduction}
\label{sec:intro}

AI-assisted software development has evolved from autocomplete-style tools such as GitHub Copilot \citep{chen2021evaluating}, through IDE-integrated assistants like Cursor~\citep{cursor2026official}, to fully agentic systems that autonomously plan multi-step modifications, execute shell commands, read and write files, and iterate on their own outputs. Claude Code~\citep{anthropic2026claudecode} is an agentic coding tool released by Anthropic~\citep{anthropic2026github}. Its official documentation describes an ``agentic loop'' that plans and executes actions toward accomplishing a goal and can call tools, evaluate results, and continue until the task is done~\footnote{\url{https://code.claude.com/docs/en/how-claude-code-works}.}. This shift from suggestion to autonomous action introduces architectural requirements that have no counterpart in completion-based tools. These requirements define a design space, a set of recurring questions spanning topics such as safety, context management, extensibility, and delegation that every coding agent must navigate. This study uses source-level analysis of Claude Code to show how one production system answers these questions.

Despite growing adoption, Anthropic publishes user-facing documentation for Claude Code but not detailed architectural descriptions. This study uses source code analysis to describe architectural design decisions. Anthropic's internal survey of 132 engineers and researchers~\citep{anthropic2025internal} reports that about 27\% of Claude Code-assisted tasks were work that would not have been attempted without the tool, suggesting that the architecture enables qualitatively new workflows rather than merely accelerating existing ones.


In this work, we first identify five human values/philosophies and thirteen design principles that motivate the architecture (\Cref{sec:values}), then organize the analysis in three parts:

\begin{enumerate}
  \item \textbf{Design-space analysis.} We identify recurring design questions (where reasoning lives, how the iteration loop is structured, what safety posture to adopt, how the extension surface is partitioned, how context is managed, how work is delegated across subagents, and how sessions persist) and analyze Claude Code's answers through a 7-component high-level structure and a 5-layer subsystem architecture, tracing each choice to specific source files (\Cref{sec:arch}). The analysis aims to build a deep understanding of the system mechanism, with the goal of informing the design of better and more powerful agent systems.
  \item \textbf{Architectural contrast with OpenClaw and Hermes Agent.}
  Beyond analyzing Claude Code itself, we also compare its design philosophy with that of two open-source agent systems, OpenClaw~\citep{openclaw2026} (a multi-channel personal assistant gateway) and Hermes Agent~\citep{hermesagent2026} (a single-process, multi-surface assistant), across six design dimensions to show how the same recurring questions produce different answers under different deployment contexts (\Cref{sec:compare}), in order to highlight both the common principles and the key differences between commercial and open-source software. This comparison helps reveal how deployment setting, product goals, safety requirements, and user assumptions shape architectural choices in different ways. By examining where these systems converge and where they diverge, our study aims to provide useful guidance and practical insights for the design of future, more capable agent systems.

  \item \textbf{Open directions for future agent systems.} Building on the design-space analysis and the OpenClaw and Hermes contrasts, \Cref{sec:future} identifies six open directions spanning the observability-evaluation gap, cross-session persistence, harness boundary evolution, horizon scaling, governance, and long-term developer capability, each drawing on empirical, architectural, and policy literature. The long-term capability question also reveals an open concern: while the Claude Code agent system amplifies the short-term capabilities of programmers and end users, it offers limited mechanisms that explicitly support long-term human improvement, deeper understanding, and sustained codebase coherence. 

\end{enumerate}


The core agent loop is a while-true cycle with state management. The surrounding subsystems for safety, extensibility, context management, delegation, and persistence make up the bulk of the implementation. Source-level analysis\footnote{Our analysis is grounded primarily in the source code, supplemented by official Anthropic documentation and selected community analysis. The appendix details the evidence base and methodology.} allows us to identify design choices, subsystem boundaries, and implementation trade-offs directly from the system itself rather than inferring them solely from product descriptions.

\paragraph{Running example.} 
To keep the architecture concrete, we trace the task ``Fix the failing test in \texttt{auth.test.ts}'' through \Cref{sec:arch,sec:turn,sec:auth,sec:ext,sec:context,sec:subagent,sec:persist}. This example illustrates how a seemingly simple user request activates multiple architectural layers, including tool invocation, permission checks, context selection, iterative repair, delegation, and session persistence.

\paragraph{Paper organization.}
\Cref{sec:values} identifies the human values and design principles that motivate the architecture. \Cref{sec:arch} introduces the high-level architecture and the design questions it answers. \Cref{sec:turn,sec:auth,sec:ext,sec:context,sec:subagent,sec:persist} each analyze a major subsystem's design choices. \Cref{sec:compare} contrasts the analysis with OpenClaw and Hermes Agent, \Cref{sec:related} positions the work against prior agent and software-engineering literature, \Cref{sec:discuss} provides discussion, and \Cref{sec:future} surveys open questions for future agent systems. \Cref{sec:conclude} concludes. The appendix describes the evidence base and methodology, package-structure notes, community reimplementations, and a companion map from new agent-system signals to the paper's design-space questions.
